\chapter*{Biodata \& End Note}
\addcontentsline{toc}{chapter}{Biodata \& End Note}

\section*{Author Biodata}

\textbf{Venkata Dhanush Kakarlamudi} (Registration Number: 23BCE8172) is an undergraduate student pursuing a Bachelor of Technology in Computer Science and Engineering at the School of Computer Science and Engineering (SCOPE), VIT-AP University, Amaravati. His academic and research interests include distributed systems reliability, statistical machine learning, streaming telemetry architectures, and Linux kernel instrumentation.

\vspace{0.4cm}

\noindent\textbf{Veera Nidhish Gondimalla} (Registration Number: 23BCE7421) is an undergraduate student pursuing a Bachelor of Technology in Computer Science and Engineering at the School of Computer Science and Engineering (SCOPE), VIT-AP University, Amaravati. His primary focus areas span high-performance computing, deep representation learning, unsupervised anomaly detection algorithms, and hardware-accelerated scientific computing.

\vspace{0.4cm}

\noindent\textbf{Settipalli Raja Rushi Praneeth Reddy} (Registration Number: 23BCE7425) is an undergraduate student pursuing a Bachelor of Technology in Computer Science and Engineering at the School of Computer Science and Engineering (SCOPE), VIT-AP University, Amaravati. His technical focus encompasses cloud infrastructure operations, site reliability engineering (SRE), machine learning pipeline automation, and predictive maintenance telemetry systems.

\vspace{0.4cm}

\noindent\textbf{Evvala Venkata Lakshmi Narasimha Veereswar} (Registration Number: 23BCE7439) is an undergraduate student pursuing a Bachelor of Technology in Computer Science and Engineering at the School of Computer Science and Engineering (SCOPE), VIT-AP University, Amaravati. His academic and research interests include scalable data engineering, applied neural networks, automated telemetry ingestion pipelines, and fault-tolerant system software.

\vspace{1.0cm}

\section*{Project End Note}

This capstone report represents the culmination of our undergraduate research efforts in large-scale system log anomaly detection. Over the course of this investigation, our focus remained steadfastly centered on empirical integrity: we audited published baselines, discovered critical undocumented preprocessing requirements, mathematically diagnosed reproducibility gaps, dismantled artificially inflated legacy benchmarks, and engineered a completely leak-free, ultra-lightweight deep learning architecture validated under realistic operating base rates.

The models, code artifacts, and experimental protocols detailed in this document demonstrate that production-grade anomaly detection does not require massive, multi-gigabyte foundation models or high-power GPU data centers. By pairing continuous semantic word representations with compact neural autoencoders and Extreme Value Theory thresholding, site reliability engineers can deploy real-time, deterministic, and microsecond-level telemetry monitoring directly on commodity CPU cores and embedded edge devices.

We hope that the empirical methodologies, the \textit{Dataset-Dependency Rule}, and the transparent disclosure of failure modes documented in this work serve as a valuable reference and foundation for future researchers and engineers advancing the reliability of autonomous systems.

